\BourbakiExerciseGroup

\begin{enumerate}
\def\labelenumi{\arabic{enumi})}
\item
  证明：若 A 是整环，则多项式环 \(A[X_{\iota}]_{\iota \in I}\) 不是主理想整环，只要以下条件之一成立：1） \(\operatorname{Card}(I) \geqslant 2\) ；2） A 不是域（考虑理想 \((\mathbf{X}_{\alpha}) + (\mathbf{X}_{\beta})\) ，对于 \(\alpha \neq \beta\) ，以及 \((a) + (\mathbf{X}_{\alpha})\) ，其中 \(a \neq 0\) 是A的一个不可逆元素）。
\item
  证明在环 \(\mathbf{K}[[\mathbf{X}]]\) ——即域 \(\mathbf{K}\) 上的形式幂级数环——中，每个不可约元素都与 \(\mathbf{X}\) 相伴。
\item
  设 \(\mathbf{A}\) 是一个整环，其中每个非空理想集合都有极大元，并且每个极大理想都是主理想；证明 \(\mathbf{A}\) 是主理想整环。（注意，如果 \(\mathfrak{a} \neq \mathbf{A}\) 是 \(\mathbf{A}\) 的理想，则存在一个极大理想 \((p) \supset \mathfrak{a}\) 使得 \(\frac{1}{p} \mathfrak{a}\) ——它是 \(\mathbf{A}\) 的理想——严格包含 \(\mathfrak{a}\) .)
\end{enumerate}

\begin{enumerate}
\def\labelenumi{\alph{enumi})}
\setcounter{enumi}{1}
\tightlist
\item
  设 K 为一个域，设 \(\mathrm{K}_{1}=\mathrm{K}((\mathrm{X}))\) 为 K 上单未定元 X 的形式幂级数域（IV，第 38 页），并设 \(B=K_{1}[Y]\) 为系数在 \(K_{1}\) 中的、关于 Y 的形式幂级数环。设 A 为 B 的由那些幂级数 \(\sum_{n=0}^{\infty}a_{n}(\mathrm{X})\mathrm{Y}^{n}\) 组成的子环，其中幂级数 \(a_{0}(\mathrm{X})\) 仅包含指数 \(\geqslant 0\) 的 X 的幂。证明主理想 (X) 是 A 的唯一极大理想，但由元素 \(YX^{-n}\) ( \(n \geqslant 0\) 为任意整数）生成的 A 的理想不是主理想。
\end{enumerate}

\begin{enumerate}
\def\labelenumi{\arabic{enumi})}
\setcounter{enumi}{3}
\tightlist
\item
  设 A 是一个整环，并设 S 是 A 的一个不包含 0 的乘法封闭子集。设 \(S^{-1}A\) 表示由元素 \(s^{-1}a\) ( \(s \in S\) , \(a \in A\) ) 在 A 的分式域 K 中组成的环（I，第 116 页）。
\end{enumerate}

\begin{enumerate}
\def\labelenumi{\alph{enumi})}
\item
  证明如果 \(\mathbf{A}\) 是主理想整环，则 \(\mathbf{S}^{-1}\mathbf{A}\) 也是主理想整环.
\item
  设 A 为主理想整环，p 为 A 的不可约元素；证明理想 \((p)\) 的补集 S 是乘法封闭的；在环 \(S^{-1}A\) 中，理想 \(pS^{-1}A\) 是唯一的极大理想，且商域 \(S^{-1}A/pS^{-1}A\) 同构于 \(A/(p)\) .
\end{enumerate}

\begin{enumerate}
\def\labelenumi{\arabic{enumi})}
\setcounter{enumi}{4}
\tightlist
\item
  设 \(\mathbf{A}\) 是一个交换环，其中每个理想都是有限生成的。
\end{enumerate}

\begin{enumerate}
\def\labelenumi{\alph{enumi})}
\item
  元素 \(p \neq 0\) （\(\mathbf{A}\) 中的元素）称为不可除的，如果它不可逆，并且关系 \(p = xy (x, y \in \mathbf{A})\) 蕴含 \(x\) 或 \(y\) 属于 \(A p\) 。证明每个元素 \(a \neq 0\) （\(\mathbf{A}\) 中的元素）都可以（至少以一种方式）写成一些不可除元与一个可逆元素的乘积（利用第七章第4页的引理1）。
\item
  交换环A中的幂等元e被称为不可分解的，如果不存在一对非零幂等元f、g使得 \(f + g = e\) 且fg = 0。两个不同的不可分解幂等元的乘积为零。证明环A中的每个幂等元都是不可分解幂等元之和（证明若 \(e\) 是幂等元，则 \(1 - e\) 也是，并利用第七章第4页的引理1）。由此推出 \(\mathbf{A}\) 是有限个环 \(\mathbf{A}_i\) 的直积，其中每个环的每个理想都是有限生成的，且单位元是唯一的非零幂等元。假设这个族至少有两个元素。那么元素 \(a = \sum_{i} a_i\) （其中 \(a_i \in \mathbf{A}_i\) 对所有 \(i\) ) 是不可除的，当且仅当存在一个指标 \(k\) 使得
\end{enumerate}

该 \(a_{k}\) 在 \(\mathbf{A}_{k}\) 中是不可除的或为零，且 \(a_{i}\) 在 \(\mathbf{A}_{i}\) 中可逆（对所有 \(i \neq k\) ）；如果 \(a_{k} = 0\) ，则 \(\mathbf{A}_{k}\) 必然是一个整环。

\begin{enumerate}
\def\labelenumi{\alph{enumi})}
\setcounter{enumi}{2}
\item
  证明：若 1 是 \(\mathbf{A}\) 的唯一非零幂等元，且若 \(x \in \mathbf{A}\) 使得 \(x^k \in \mathbf{A}x^{k+1}\) 对于某个整数 \(k \geqslant 1\) 成立，则要么 \(x\) 可逆，要么 \(x^k = 0\) （若 \(x^k = ax^{k+1}\) ，则考虑 \(a^k x^k\) ).
\item
  设 \(p \neq 0\) 是 A 中的零因子。证明存在一个有限序列 \((a_k)_{1 \leqslant k \leqslant m}\) 由 A 的非零元素组成，使得 \(0 = pa_1\) , \(a_k = pa_{k+1}\) 对于 \(1 \leqslant k < m\) , \(a_m \notin \mathbf{A}_p\) ，以及 \((a_k) \neq (a_{k+1})\) 对于 \(1 \leqslant k < m\) .
\end{enumerate}

\begin{enumerate}
\def\labelenumi{\arabic{enumi})}
\setcounter{enumi}{5}
\tightlist
\item
  一个交换环 A 被称为主理想环，如果 A 的每个理想都是主理想。a) 证明主理想环的商环以及主理想环的有限乘积都是主理想环。
\end{enumerate}

\begin{enumerate}
\def\labelenumi{\alph{enumi})}
\setcounter{enumi}{1}
\item
  证明每个主理想环 A 都是某个有限族 \((\mathbf{A}_{i})_{1\leq i\leq m}\) 的直积，族中每个环都是主理想环，且其单位元是唯一的非零幂等元（利用习题5，b））。
\item
  设 A 是一个主理想环，其中单位元是唯一的非零幂等元，并且 A 含有除 0 以外的零因子。证明 A 中存在一个不可除的幂零元 p。（利用习题 5 a），证明存在一个不可除元 \(p \neq 0\) 属于 A，它是一个零因子；应用习题5 d)，然后证明主理想 Ab （\(y \in A\) 满足 \(yp^{m} \in Ap^{m+1}\) ）就是整个 A；最后应用习题5，c)。d) 在与c)相同的假设和记号下，假设 \(p^{m} = 0\) 且 \(p^{m-1} \neq 0\) 。证明 \(p^{m-1}\) 的零化子是 Ap（否则它将是一个主理想 Ac，满足 \(p \in Ac\) 且 \(c \notin Ap\) ；利用 p 不可除这一事实，推出 c 将是可逆的）；由此推出，对于 \(1 \leq k \leq m-1\) ， \(p^{m-k}\) 的零化子是 \(Ap^{k}\) 。由此得出理想 Ap 是极大的（仿照 d) 中第一个断言的论证，利用 1 - xp 在 A 中可逆，对所有 \(x \in A\) .
\item
  由 \(d\) ) 推出，对所有 \(x \in \mathbf{A}\) ，存在一个整数 \(k \geqslant 0\) 以及一个可逆元素 \(u\) 使得 \(x = up^k\) ，并且 \(\mathbf{A}\) 的理想只有诸 \(Ap^k\) ( \(0 \leqslant k \leqslant m\) ) （观察到 \(\mathbf{A}\) 中每个不属于 \(Ap\) 的元素都在 \(\mathbf{A}\) 中可逆).
\item
  设 A 为主理想环，令 \((x_{i})\) 为 A 的一族元素，并令 \(Ad\) 为由族 \((x_{i})\) 生成的理想；证明存在元素 \(x_{i}^{\prime}\) 属于 A，使得 \(x_{i} = dx_{i}^{\prime}\) 对所有 \(i\) 成立，以及 \(\sum_{i} Ax_{i}^{\prime} = A\) （化归到 1 是唯一的
\end{enumerate}

非零幂等元的情形）。

\(g)\) 给出一个主理想环和一个不可除元素 \(p\) （ \(\mathbf{A}\) 中的）的例子，使得理想 \(Ap\) 不是极大的（利用习题5， \(b\) )).

\begin{enumerate}
\def\labelenumi{\arabic{enumi})}
\setcounter{enumi}{6}
\tightlist
\item
  给定一个整环 \(\mathbf{A}\) ，我们称一个映射 \(w\) ——从 \(\mathbf{A}' = \mathbf{A} - \{0\}\) 到 \(\mathbf{N}\) 的映射——为欧几里得函数，如果它满足以下条件：
\end{enumerate}

\((S_{1}) w(xy) \geqslant w(y)\) 对于每一对 \(x, y\) ，其元素属于 \(\mathbf{A}'\) .

（S \(_{II}\) ) 如果 \(a \in \mathbf{A}'\) 且 \(b \in \mathbf{A}'\) ，则存在元素 \(q\) , \(r\) 属于 \(\mathbf{A}\) ，使得 \(a = bq + r\) ，并且要么 \(r = 0\) ，要么 \(w(r) < w(b)\) .

如果 A 上存在一个欧几里得函数，则称 A 为欧几里得整环。

\begin{enumerate}
\def\labelenumi{\alph{enumi})}
\item
  证明： \(\mathbf{Z}\) 和 \(\mathbf{K}[\mathbf{X}]\) （其中 \(\mathbf{K}\) 是域）都是欧几里得整环。
\item
  证明每个欧几里得整环都是主理想整环（如果 \(\mathfrak{a}\) 是 \(\mathbf{A}\) 的理想，选择 \(a \in \mathfrak{a}\) 使得 \(w(a)\) 尽可能小）。
\item
  证明高斯整数环 A（VII，第～1页）是关于函数 \(w(a + bi) = a^2 + b^2 = \mathbf{N}_{\mathbf{Q}(i) / \mathbf{Q}}(a + bi)\) 的欧几里得整环（注意，对所有 \(x \in \mathbf{Q}(i)\) ，存在 \(y \in \mathbf{A}\) 使得 \(\mathbf{N}(x - y) \leqslant 1/2\) ).
\end{enumerate}

\(d\) ）证明二次扩张环 \(\mathbf{B}\) （\(\mathbf{Z}\) 的、以 \((1,e)(e^2 = 2)\) 为基的环）是关于函数

\[
w (a + b e) = \left| a ^ {2} - 2 b ^ {2} \right| = \left| N _ {Q (e) / Q} (a + b e) \right|
\]

（对所有 \(x \in \mathbf{Q}(e)\) ，存在 \(y \in B\) 使得 \(|\mathbf{N}(x - y)| \leqslant 1/2\) )的欧几里得整环.

\begin{enumerate}
\def\labelenumi{\arabic{enumi})}
\setcounter{enumi}{7}
\tightlist
\item
  设 \(\mathbf{K}\) 是一个二次域，即 \(\mathbf{Q}\) 的次数为 2 的扩张。
\end{enumerate}

\begin{enumerate}
\def\labelenumi{\alph{enumi})}
\item
  证明： \(\mathbf{K} = \mathbf{Q}(e)\) ，其中 \(e^2\) 是一个有理整数 \(d \neq 1\) ，不被 \(>1\) 的平方数整除（在 \(\mathbf{Z}\) 中）.
\item
  如果 \(\alpha = a + be\) ( \(a, b \in \mathbf{Q}\) ），令 \(\overline{\alpha}\) 表示共轭 \(a - be\) （即 \(\alpha\) 的共轭）。我们称 \(\alpha\) 是二次域 \(\mathbf{K}\) 的整数，如果其范数 \(\alpha\overline{\alpha}\) 及其迹 \(\alpha + \overline{\alpha}\) 是有理整数。证明这些整数（\(\mathbf{K}\) 的整数）构成一个环 \(\mathbf{A}\) （注意如果 \(\alpha\) 和 \(\beta\) 是 \(\mathbf{K}\) 的整数，那么有理数 \(\alpha\beta + \overline{\alpha\beta}\) 和 \(\alpha\overline{\beta} + \overline{\alpha}\beta\) 的和与积是整数，因此这些有理数是 \(\mathbf{Z}\) 上某个首一多项式的根；由此推断它们是有理整数）。
\item
  证明 \(\mathbf{K} = \mathbf{Q}(e)\) 的整数构成一个 \(\mathbf{Z}\) -模，其基为 \((1, e)\) ，当 \(d - 1 \not\equiv 0 \pmod{4}\) 时；或为 \((1/2(1 + e), 1/2(1 - e))\) ，当 \(d - 1 \equiv 0 \pmod{4}\) 时（若 \(a + be\) 是 \(\mathbf{K}\) 的整数，验证 \(2a\) 和 \(2b\) 是满足 \((2a)^2 - d(2b)^2 \equiv 0(4)\) 的有理整数，并研究它们是否可能为奇数）。证明这些基的判别式（III，第～549页，定义～3）分别为 \(4d\) 和 \(d\) 。
\end{enumerate}

\begin{enumerate}
\def\labelenumi{\arabic{enumi})}
\setcounter{enumi}{8}
\item
  *（闵可夫斯基定理）若 \(a, b, c, d\) 是实数，且 \(\mathbf{D} = \begin{vmatrix} a & b \\ c & d \end{vmatrix}\) ，则不等式 \(|na + mb| \leqslant \mathbf{A}\) , \(|nc + md| \leqslant \mathbf{B}\) 有非零整数解 \((n, m)\) ，只要 \(\mathbf{A} > 0\) , \(\mathbf{B} > 0\) , \(\mathbf{D} > 0\) 且 \(\mathbf{AB} \geqslant \mathbf{D}\) ，且这样的解只有有限多个（考虑离散子群 \(\mathbf{G}\) ——\(\mathbf{R}^2\) 中由 \((a, c)\) 和 \((b, d)\) 生成的子群——，并证明矩形 \(\mathbf{E}\) ——由关系 \(0 \leqslant x \leqslant \mathbf{A}\) , \(0 \leqslant y \leqslant \mathbf{B}\) 定义，其面积大于在向量 \((a, c)\) 和 \((b, d)\) 上构造的平行四边形的面积——包含两个模 \(\mathbf{G}\) 同余的不同向量；用反证法论证，注意到否则 \(p^2\) 个矩形——它们由 \(\mathbf{E}\) 借助平移（ \(ha + kb\) , \(hc + kd\) ）（其中 \(0 \leqslant h < p\) 且 \(0 \leqslant k < p\) ）得到——将是两两不相交的）。
\item
  设 \(d\) 为一个不被 \(>1\) 的平方数整除的有理整数。在二次域 \(\mathbf{K} = \mathbf{Q}(\sqrt{d})\) 中，我们将考察单位，即环 \(\mathbf{A}\) ——\(\mathbf{K}\) 的整数环——中的可逆元素（习题8，b））。
\end{enumerate}

\begin{enumerate}
\def\labelenumi{\alph{enumi})}
\item
  如果 \(d < 0\) （“虚二次域”），则证明 \(K\) 的单位只有 1 和 -1，但 \(d = -1\) 例外，此时 \(i\) 和 -i 也是单位；以及 \(d = -3\) 的情形，此时四个数 \(1/2(\pm 1 \pm \sqrt{-3})\) 也是单位。
\item
  从现在起假设 \(d > 0\) （“实二次域”），并假设 \(\mathbf{K}\) 嵌入到 \(\mathbf{Q}\) 的一个极大序扩张中 (VI, 第～25页)，* 例如在 \(\mathbf{R}\) 中. * 证明正单位构成 K 的单位群 U 的子群 \(U'\) ，并且 U 是 \(U'\) 与 \(\{-1,1\}\) 的直积.
\item
  证明 K 的单位是范数为 1 或 -1 的整数。
\end{enumerate}

\(d\) ) 证明群 \(\mathbf{U}^{\prime}\) 是非平凡的（令 D 表示整数基的判别式（习题 8， \(c\) )；证明 K 中的整数 \(\alpha\) （满足 \(\mathbf{N}(\alpha)\leqslant \mathbf{D}\) ）被划分为模 U 的有限多个同余类；在每个类中选取一个整数 \(\beta_{i}\) ；如果 B 是一个有理数，使得 \(0 < \mathrm{B} < \inf (|\beta_i|)\) ，证明（习题9）存在 K 的一个整数 \(\beta\) ，满足 \(0 < |\beta |\leqslant \mathrm{B}\) 且 \(\mathbf{N}(\beta)\leqslant \mathbf{D}\) ；由此推出某个 \(\beta \beta_{i}^{-1}\) 是一个单位 \(\eta\) ，满足 \(|\eta| < 1\) ).

\begin{enumerate}
\def\labelenumi{\alph{enumi})}
\setcounter{enumi}{4}
\item
  证明： \(\mathbf{U}'\) 同构于 \(\mathbf{Z}\) （证明存在一个单位 \(\varepsilon\) ，使得 \(|\varepsilon| < 1\) 且 \(|\varepsilon|\) 在具有此性质的数中尽可能大，利用习题9以及事实 \(|\overline{\xi}| = |\xi|^{-1}\) （对单位成立）。
\item
  证明在 \(\mathbf{Q}(\sqrt{2})\) 中，元素 \(\sqrt{2} - 1\) 是 \(\mathbf{U}'\) 的一个生成元，且其范数为 \(-1\) .
\end{enumerate}

\(g\) ) 证明如果 \(d\) 有一个素因子 \(p\) （形如 \(4n - 1\) ），则 \(\mathbf{U}'\) 的每个生成元都具有范数 1。（注意， \(-1\) 在域 \(\mathbf{F}_p\) 中不能是平方数，利用 V，第 93 页，命题 1，c)。

¶ 11) 我们研究这样的二次域 \(\mathbf{Q}(\sqrt{d})\) ：映射 \(\alpha \mapsto |\mathrm{N}(\alpha)|\) 是该域整数上的欧几里得函数（习题 7）。

\begin{enumerate}
\def\labelenumi{\alph{enumi})}
\item
  证明其充分必要条件为：对所有 \(\alpha \in \mathbf{Q}(\sqrt{d})\) ，存在该域的一个整数 \(\beta\) 使得 \(|\mathrm{N}(\alpha - \beta)| < 1\) ；写出此条件，将 \(\alpha\) 和 \(\beta\) 用该域的整数基表示（习题8，c））。
\item
  对于虚二次域的情形 \((d < 0)\) ，证明使 \(|\mathbf{N}(\alpha)|\) 是欧几里得函数的 d 的值只有 -1、-2、-3、-7、-11。
\item
  对于实二次域，我们将仅限于 \(d \equiv 2\) 或 3（模 4）的情形。证明此时仅有有限多个 \(d\) 的值使得 \(|\mathrm{N}(\alpha)|\) 是欧几里得函数。（若 \(|\mathrm{N}(\alpha)|\) 是欧几里得函数，且 \(a\) 是一个自然数，使得 \(0 \leqslant a \leqslant d\) 且 \(a \equiv t^2(d)\) ，则证明 \(a\) 或 \(d - a\) 之一具有形式 \(x^2 - dy^2\) ，其中 \(x\) 和 \(y\) 是有理整数，这通过考虑元素 \(t / \sqrt{d}\) （属于 \(\mathbf{Q}(\sqrt{d})\) ）来证明；接下来证明，对于足够大的 \(d\) ，存在一个奇数 \(z\) 使得 \(5d < z^2 < 6d\) （若 \(d \equiv 3 \mod 4\) ）或 \(2d < z^2 < 3d\) （若 \(d \equiv 2 \mod 4\) ）；然后取 \(a = z^2 - d[z^2/d]\) ，并证明在这些条件下，\(a\) 与 \(d - a\) 都不具有形式 \(x^2 - dy^2\) （通过考察模 8 的剩余类）。证明 \(d\) 的值，当这样的整数 \(z\) 不存在时，只有 3、7、11、19、35、47 和 59（对于 \(d \equiv 3 \mod 4\) ）以及 2、6、14 和 26（对于 \(d \equiv 2 \mod 4\) ）。证明对于 \(d = 2\) 和 \(d = 3\) 确实存在欧几里得函数。
\end{enumerate}

\begin{enumerate}
\def\labelenumi{\arabic{enumi})}
\setcounter{enumi}{11}
\tightlist
\item
  证明二次域 \(\mathbf{Q}(\alpha)\) （其中 \(\alpha^2 = -5\) ）的整数环不是主理想整环（参见第六章，第34页，习题27）。对 \(\mathbf{Q}(\beta)\) （其中 \(\beta^2 = 10\) ）提出同样的问题（我们有
\end{enumerate}

\[
2. 3 = (4 + \beta) (4 - \beta)).
\]

\begin{enumerate}
\def\labelenumi{\arabic{enumi})}
\setcounter{enumi}{12}
\item
  设 \(f\) 是 \(\mathbf{Z}\) 上一个未定元的非常数多项式。证明对于每个整数 \(k \geqslant 1\) ，存在无穷多个整数 \(n \in \mathbf{N}\) 使得 \(f(n)\) 非零且至少被 \(k\) 个不同的素数整除（可以对 \(k\) 用归纳法：如果 \(n \in \mathbf{N}\) 使得 \(f(n)\) 非零且至少被 \(k\) 个不同的素数整除，则考虑多项式 \(f(n + Xf(n)^2)\) .
\item
  形如 \(4n - 1\) （相应地 \(6n - 1\) ）的素数集合是无限的（方法与命题 5（VII，第～5 页）相同；注意除 2（相应地，除 2 和 3）以外的每个素数都具有形式 \(4n \pm 1\) （相应地 \(6n \pm 1\) )).
\item
  形如 \(2^k + 1\) ( \(k \geqslant 1\) 为整数）的数只有当 \(k\) 是 2 的幂时才可能是素数。如果素数 \(p\) 整除 \(2^{2^n} + 1\) ，那么它不能整除 \(2^{2^m} + 1\) ，对于 \(m > n\) ( \(2^{2^m} \equiv 1(p)\) ）；利用这一点给出命题5的一个新证明。证明 \(2^{2^5} + 1\) 不是素数（令 \(a = 2^7\) 和 \(b = 5\) ；证明 \(1 + ab - b^4 = 2^4\) 且 \(1 + ab\) 整除
\end{enumerate}

\[
n = 2 ^ {4} a ^ {4} + 1 = (1 + a b - b ^ {4}) a ^ {4} + 1.
\]

\begin{enumerate}
\def\labelenumi{\arabic{enumi})}
\setcounter{enumi}{15}
\item
  形如 \(a^k - 1\) ( \(a\) 和 \(k\) 为整数， \(a > 0\) , \(k > 1\) ) 的数只有当 \(a = 2\) 且 \(k\) 是素数时才可能是素数。
\item
  证明如果 \(a^n + 1\) 是素数，其中 \(a > 1\) 且 \(n > 1\) ，则 \(a\) 是偶数且 \(n = 2^m\) .
\item
  设 M 是由有限个整数 \(q_{i} > 1\) ( \(1 \leqslant i \leqslant m\) ) 生成的乘法幺半群。证明：对于每个整数 k \textgreater{} 0，存在 k 个连续整数，它们都不属于 M（如果 \(h_{1}, \ldots, h_{m}\) 是 m 个 \textgreater{} 1 的整数，且 k 是都不被任何 \(h_{i}\) 整除的连续整数的最大个数，则 \((k + 1)^{-1} \leqslant h_{1}^{-1} + \cdots + h_{m}^{-1}\) ；现在注意到，对所有 r \textgreater{} 0，M 中每个足够大的数都至少被一个 \(q_{i}^{+}\) 整除）。由此结果给出 VII 第5页命题5的一个新证明。
\item
  对每个整数 \(n > 0\) ，证明素数 \(p\) 在 \(n!\) 的素因子分解中的指数是 \(\sum_{k=1}^{\infty} [np^{-k}]\) （其中只有有限多个非零项）。
\item
  * 设 \(\pi(x)\) 为小于实数 \(x > 0\) 的素数的个数。对于每个整数 \(n > 0\) ，证明
\end{enumerate}

\[
n ^ {\pi (2 n) - \pi (n)} \leqslant \binom {2 n} {n} \leqslant (2 n) ^ {\pi (2 n)}
\]

（注意 \(\binom{2n}{n}\) 是所有素数 \(q\) （满足 \(n < q \leqslant 2n\) ）的乘积的倍数）；另一方面，利用习题19证明 \(\binom{2n}{n}\) 整除乘积 \(\prod p^{r(p)}\) ，其中 \(p\) 遍历素数集合 \(\leqslant 2n\) ，而 \(r(p)\) 是满足 \(p^{r(p)} \leqslant 2n\) 的最大整数。由此结果推出存在两个实数 \(a\) 和 \(b\) 使得 \(0 < a < b\) 且

\[
a n (\log n) ^ {- 1} \leqslant \pi (n) \leqslant b n (\log n) ^ {- 1}
\]

（注意 \(2^{2n}(2n + 1)^{-1} \leqslant (2n)! (n!)^{-2} \leqslant 2^{2n}\) ). *

\(^{1}\) 已证明 \(\pi(x) \sim \frac{x}{\log x}\) 当 x 趋于 \(+\infty\) 时（参见例如 A. E. INGHAM，《素数的分布》（剑桥丛刊，第 30 号），剑桥大学出版社，1932 年）。

\begin{enumerate}
\def\labelenumi{\arabic{enumi})}
\setcounter{enumi}{20}
\tightlist
\item
  证明对于 \(m \geqslant 1\) 和 \(n \geqslant 1\) ，有理数 \(\sum_{k=n}^{n+m} k^{-1}\) 绝不是整数（若
\end{enumerate}

\(2^{q}\) 是整除数 \(n, n+1, \ldots, n+m\) 中至少一个的 2 的最大幂，则它只整除这些数中的一个）。

\begin{enumerate}
\def\labelenumi{\arabic{enumi})}
\setcounter{enumi}{21}
\item
  设 \(a\) 和 \(b\) 是两个 \(> 0\) 的整数，设 \(d\) 是 \(a\) 和 \(b\) 的最大公因数，并设 \(n\) 是一个 \(\geqslant 1\) 的整数；令 \(a = da_1\) 和 \(b = db_1\) ；证明分数 \(a(a + b) \ldots (a + (n - 1)b) / n!\) 的既约表达式中的分母只具有整除 \(b_1\) 且 \(\leqslant n\) 的素因子（注意，如果 \(p\) 是一个不整除 \(b_1\) 的素数，则 \(b_1\) 模 \(p^r\) 可逆，对所有 \(r > 0\) ). 由此直接证明二项式系数 \(n! / (m!(n - m)! )\) 是整数。
\item
  设 \(n > 1\) 是一个有理整数；证明这 \(n - 1\) 个二项式系数 \(\binom{n}{k}\) ( \(1 \leqslant k \leqslant n - 1\) ) 的最大公因数等于 1，如果 \(n\) 被两个不同的素数整除；等于 \(p\) ，如果 \(n\) 是素数 \(p\) 的幂。
\item
  \begin{enumerate}
  \def\labelenumii{\alph{enumii})}
  \tightlist
  \item
    设 \(n = \prod_{k} p_k^{\nu_k} > 0\) 是一个用其素因子表示的有理整数。
  \end{enumerate}
\end{enumerate}

证明 n 的正因数之和由公式给出

\[
\sigma (n) = \prod_ {k} \frac {p _ {k} ^ {\nu_ {k} + 1} - 1}{p _ {k} - 1}.
\]

\begin{enumerate}
\def\labelenumi{\alph{enumi})}
\setcounter{enumi}{1}
\tightlist
\item
  我们说 \(n\) 是完全数，如果 \(2n = \sigma(n)\) 。证明一个偶数 \(n\) 是完全数，当且仅当它具有形式 \(2^h(2^{h+1} - 1)\) ，其中 \(h \geqslant 1\) 且 \(2^{h+1} - 1\) 是素数。（如果 \(n\) 具有形式 \(2^m q\) ，其中 \(m \geqslant 1\) 且 \(q\) 是奇数，证明 \(q\) 能被 \(2^{m+1} - 1\) 整除，并由此推出，如果 \(q\) 不是素数，那么我们将有 \(\sigma(n) > 2n\)）。
\end{enumerate}

\begin{enumerate}
\def\labelenumi{\arabic{enumi})}
\setcounter{enumi}{24}
\item
  设 \(a\) 和 \(b\) 是两个 \(> 0\) 的互素整数。若 \(a\) 和 \(b\) 都 \(> 2\) ，则存在整数 \(x\) 和 \(y\) 使得 \(ax + by = 1\) 且 \(|x| < 1/2b\) , \(|y| < 1/2a\) ；此外，整数 \(x\) 和 \(y\) 由这些条件唯一确定。当 \(a, b\) 之一等于 2 时情况如何？
\item
  设 \(x, y\) 是两个互素的整数；则每个整数 \(z\) （满足 \(|z| < |xy|\) ）都可以以一种或两种方式表示为 \(z = ux + vy\) 的形式，其中 \(u\) 和 \(v\) 是满足 \(|u| < |y|\) 和 \(|v| < |x|\) 的整数。若 \(z\) 能以这种形式用两种不同方式表示，则 \(x\) 和 \(y\) 都不能整除 \(z\) ；例子 \(x = 5, y = 7, z = 12\) 表明其逆命题不成立。
\item
  设 \(f\) 和 \(g\) 是 \(\mathbf{K}[\mathbf{X}]\) 中的两个互素多项式；证明每个多项式 \(h\) （满足 \(\deg(h) < \deg(fg)\) ）都可以唯一地表示为 \(h = uf + vg\) 的形式，其中 \(u\) 和 \(v\) 是满足 \(\deg(u) < \deg(g)\) 和 \(\deg(v) < \deg(f)\) 的多项式。首先给出一个与习题26类似的证明（基于欧几里得算法），然后给出一个将该问题视为线性问题的证明。
\item
  \begin{enumerate}
  \def\labelenumii{\alph{enumii})}
  \tightlist
  \item
    设K是一个代数闭的交换域，设f是K{[}X{]}中的一个非常数的首一多项式，并设g是K{[}X, Y{]}中的一个多项式；假设对于f的每一个根 \(\alpha_{i}\) ，都存在一个根 \(\beta_{i}\) ——即 \(g(\alpha_{i}, \mathbf{Y}) = 0\) 的根——使得 \(\frac{\partial g}{\partial \mathbf{Y}} (\alpha_{i}, \beta_{i}) \neq 0\) 。则证明存在一个多项式 \(h(\mathbf{X}) \in \mathbf{K}[\mathbf{X}]\) 使得 \(g(\mathbf{X}, h(\mathbf{X})) \equiv 0 \pmod{f(\mathbf{X})}\) （利用命题4，化归到情形 \(f(\mathbf{X}) = \mathbf{X}^{m}\) （对某个 \(m \geqslant 1\) ）；注意环 \(K[X]/(X^{m})\) 同构于商环 \(K[[X]]/(X^{m})\) ，并利用 IV 第 37 页的推论）。
  \end{enumerate}
\end{enumerate}

\begin{enumerate}
\def\labelenumi{\alph{enumi})}
\setcounter{enumi}{1}
\tightlist
\item
  特别地证明，若 m 是不被 K 的特征整除的整数，则对每一对互素的首一多项式 f 和 \(f_{0}\) （在 K{[}X{]} 中），存在一个多项式 \(h(\mathbf{X})\) （在 K{[}X{]} 中），使得 \((h(\mathbf{X}))^{m} \equiv f_{0}(\mathbf{X}) \pmod{f(\mathbf{X})}\) 。
\end{enumerate}

